\documentclass[14pt,a4paper]{extarticle}
\usepackage[T2A]{fontenc}
\usepackage[utf8]{inputenc}
\usepackage[english,russian]{babel}
\usepackage[left=25mm,right=20mm,top=20mm,bottom=20mm]{geometry}
\usepackage{amsmath,amssymb,booktabs,array,longtable}
\usepackage{graphicx}
\usepackage{enumitem}
\usepackage{titlesec}
\usepackage[unicode,hidelinks]{hyperref}
\setlength{\parindent}{1.25cm}
\setlength{\emergencystretch}{3em}
\setlist{nosep}
\setcounter{tocdepth}{2}
\newcommand{\Sop}{S}
\newcommand{\Cg}{C_g}
\newcommand{\dd}{\mathrm{d}}
\newcommand{\Rnum}{\mathbb{R}}
\newcommand{\code}[1]{\texttt{#1}}
\newcommand{\emptyfield}[1]
{\par\vspace{5mm}\noindent\textbf{#1}\par\vspace{9mm}}
\newcommand{\pt}[1]{\textbf{#1}\ }
% Численные значения перенесены из article/tables/numbers.tex.

\begin{document}
% Титульный лист по образцу «Отчет по проекту.docx».
% Tempora — шрифт семейства Times с кириллицей, доступный в TeX Live.
\begin{titlepage}
\newgeometry{left=25.4mm,right=25.4mm,top=25.4mm,bottom=25.4mm}
\fontfamily{Tempora-TLF}\fontsize{14}{16}\selectfont
\setlength{\parindent}{0pt}
\begin{center}
федеральное государственное автономное образовательное учреждение\\
высшего образования\\
\textbf{«Национальный исследовательский университет\\
«Высшая школа экономики»}

\vspace{6mm}
\textbf{Московский институт электроники и математики\\
им. А.\,Н. Тихонова}

\vspace{29mm}
\textbf{Отчет\\по проектной работе}

\vspace{3mm}
\textit{«Методы нахождения решений линейных дифференциальных уравнений второго порядка с переменными коэффициентами»}
\end{center}

\vspace{18mm}
\begin{flushright}
{\fontsize{12}{14}\selectfont\textbf{Выполнили студенты}}
\end{flushright}
\vspace{3mm}
{\fontsize{12}{15}\selectfont
\newcommand{\studentfield}[1]{%
\begin{minipage}[t]{0.47\textwidth}
\raggedleft
#1\par
\vspace{4mm}
\rule{45mm}{0.4pt}\\[-1mm]
\textit{(подпись)}
\end{minipage}}
\noindent
\studentfield{Авдеев Евгений Владимирович, БПМ233}\hfill
\studentfield{Евсеев Кирилл Александрович, БПМ231}\par
\vspace{8mm}
\noindent
\studentfield{Тимофеев Михаил Юрьевич, БПМ231}\par
}

\vspace{13mm}
\textbf{Руководитель проекта:}\par
\vspace{6mm}
Парусникова Анастасия Владимировна\par
\vspace{5mm}
\rule{55mm}{0.4pt}\\[-1mm]
{\fontsize{12}{14}\selectfont\textit{(подпись)}}

\vfill
\begin{center}Москва, 2026\end{center}
\end{titlepage}
\restoregeometry

\section*{Информация о составе проектной команды и контакты}
\phantomsection
\addcontentsline{toc}{section}{Информация о составе проектной команды и контакты}

\subsection*{Парусникова Анастасия Владимировна:}
\textbf{Роль:} руководитель проекта\\
\textbf{Контакты:} aparusnikova@hse.ru
\subsection*{Авдеев Евгений Владимирович}
\textbf{Роль:} научный работник-исследователь\\
\textbf{Контакты:} evavdeev@edu.hse.ru
\subsection*{Евсеев Кирилл Александрович}
\textbf{Роль:} научный работник-исследователь\\
\textbf{Контакты:} kaevseev@edu.hse.ru
\subsection*{Тимофеев Михаил Юрьевич}
\textbf{Роль:} промпт-инженер\\
\textbf{Контакты:} mitimofeev@edu.hse.ru
\clearpage

\section*{Реферат}
\phantomsection
\addcontentsline{toc}{section}{Реферат}
\subsection*{Выбор численного метода с помощью машинного обучения}
\begingroup
\titlespacing*{\subsubsection}{0pt}{5pt}{2pt}
\subsubsection*{Задача}
Рассматривается уравнение
\[
a(x)f''+b(x)f'+\bigl(c(x)-\lambda\bigr)f=-g(x).
\]
Его решение вычисляется четырьмя численными методами: методом Ремизова,
методом Монте-Карло, методом Рунге--Кутты четвёртого порядка (RK4)
и методом мультипликативного интеграла. Все методы сравниваются
по ошибке в точке $x=0$:
\[
E_m=\bigl|f_m(0)-f(0)\bigr|,
\]
где $f_m$ --- решение, найденное методом $m$, а $f$ --- точное решение.
Метод достигает точности $\varepsilon=0.05$, если $E_m\le0.05$.

\subsubsection*{Цель работы}
Оценить точность методов при фиксированных настройках и проверить,
можно ли до решения уравнения, только по его коэффициентам и правой
части, предсказать метод с меньшей ошибкой.

\subsubsection*{Методика}
\begin{itemize}
\item Построено 1000 уравнений с известным точным решением: задаются
$a,b,c,\lambda$ и $f$, а $g$ вычисляется подстановкой в уравнение.
Уравнения образуют 200 семейств по 5 уравнений.
\item Каждое уравнение описывается 258 числами (признаками): значения
функций $a$, $b$, $c$, $g$ на сетке, их статистики, границы коэффициентов
и характеристики формул коэффициентов.
\item Модель выбора между методами Ремизова и Монте-Карло --- CatBoost
(сумма 120 решающих деревьев). Она обучается на 747 уравнениях, порог
выбора подбирается на 100, проверка проводится на 150. Уравнения одного
семейства не попадают в разные части.
\item Основной критерий --- средняя избыточная ошибка: на сколько
в среднем ошибка выбранного метода больше ошибки лучшего из двух методов.
\item На тех же признаках и выборках обучена отдельная модель CatBoost
для выбора из трёх методов: Ремизова, Монте-Карло и мультипликативного
интеграла.
\end{itemize}

\subsubsection*{Результат 1. Методы Ремизова и Монте-Карло}
На 997 уравнениях (для трёх расчёт методом Ремизова не завершился)
метод Ремизова точнее в 881 случае, Монте-Карло --- в 116. При этом
у метода Ремизова встречаются редкие очень большие ошибки (до 842),
а ошибка Монте-Карло не превышает $0.57$.

\subsubsection*{Результат 2. Модель выбора из двух методов}
На 150 тестовых уравнениях модель нашла 6 из 9 уравнений, где точнее
Монте-Карло, но 14 раз выбрала Монте-Карло ошибочно; в двух случаях
из-за этого ошибка превысила $0.05$. Малая средняя ошибка модели
объясняется в основном одним уравнением, где ошибка метода Ремизова
равна 248. Без этого уравнения постоянный выбор метода Ремизова даёт
меньшую среднюю ошибку, чем модель.

\subsubsection*{Результат 3. Четыре метода}
RK4 даёт наименьшую ошибку в точке $x=0$ на всех 997 уравнениях
и достигает точности $0.05$ на всех 1000. Поэтому на этом наборе выбор
из четырёх методов тривиален, и модель для него не обучалась.

\subsubsection*{Результат 4. Выбор из трёх методов}
Без RK4 лучшие методы распределены так: на 826 уравнениях точнее
остальных двух метод Ремизова, на 106 --- Монте-Карло, на 65 --- метод
мультипликативного интеграла. Модель верно определила лучший метод на 133 из 150
тестовых уравнений (88.67\%), а выбранный ею метод достиг точности $0.05$
на 147 (98\%). Однако постоянный выбор метода Ремизова верен на 136
уравнениях (90.67\%) и тоже достигает точности на 147, поэтому
преимущества у модели нет.

\subsubsection*{Выводы}
\begin{enumerate}
\item Метод Ремизова чаще точнее Монте-Карло, но изредка даёт очень
большую ошибку.
\item Выбор более точного из двух методов и достижение точности $0.05$ ---
разные свойства: верно выбранный метод может не достигать точности.
\item Модель выбора из двух методов находит часть уравнений, где точнее
Монте-Карло, но ценой многих ошибочных выборов. Ни она, ни модель выбора
из трёх методов не дают устойчивого выигрыша перед постоянным выбором
метода Ремизова.
\item На построенном наборе при выбранных настройках самым точным
в точке $x=0$ оказался RK4.
\end{enumerate}
\endgroup
\clearpage
\tableofcontents
\clearpage

\section{Основная часть}
\subsection{Введение}
\subsubsection{Объект и предмет исследования}
Объект исследования --- линейные неоднородные дифференциальные
уравнения второго порядка с переменными коэффициентами. Предмет
исследования --- точность четырёх численных методов их решения (метода
Ремизова, метода Монте-Карло, RK4 и метода мультипликативного интеграла)
и возможность заранее, по коэффициентам уравнения, выбрать метод
с меньшей ошибкой.

\subsubsection{Цель и задачи}
Цель работы --- оценить точность численных методов при фиксированных
настройках и проверить, можно ли до решения уравнения предсказать метод
с меньшей ошибкой.

Для достижения цели решались следующие задачи:
\begin{enumerate}
\item Построить набор из 1000 уравнений с известными точными решениями.
\item Вычислить ошибки четырёх методов в точке $x=0$ при фиксированных
настройках.
\item Описать каждое уравнение набором признаков, которые вычисляются
до его решения.
\item Обучить модель выбора между методами Ремизова и Монте-Карло
и сравнить её с простыми правилами выбора.
\item Разобрать все неверные выборы модели на тестовой выборке.
\item Сравнить четыре метода и определить, имеет ли смысл модель выбора
из четырёх методов.
\item Обучить модель выбора из трёх методов --- Ремизова, Монте-Карло
и мультипликативного интеграла --- и сравнить её с постоянным выбором
одного метода.
\end{enumerate}

\subsection{Ход работы}
\subsubsection{Математическая постановка}
Рассматривается уравнение
\begin{equation}\label{eq:ode}
a(x)f''(x)+b(x)f'(x)+\bigl(c(x)-\lambda\bigr)f(x)=-g(x),
\end{equation}
где коэффициенты $a,b,c$, правая часть $g$ и число $\lambda$ заданы,
а $f$ --- искомое решение.

Методы Ремизова и Монте-Карло решают уравнение~\eqref{eq:ode} на всей
прямой и начальных условий не используют. RK4 и метод мультипликативного
интеграла решают задачу Коши на отрезке $[-1,1]$ с начальными условиями
$f(-1)$ и $f'(-1)$, сохранёнными при построении уравнения. Для нового
уравнения, точное решение которого неизвестно, такие начальные условия
должны быть известны заранее.

Все методы сравниваются по ошибке в одной точке $x=0$:
\begin{equation}\label{eq:ml-error}
E_m=\bigl|f_m(0)-f(0)\bigr|,
\end{equation}
где $m$ --- метод, $f_m$ --- найденное им решение, $f$ --- точное
решение. Будем говорить, что метод достигает точности $\varepsilon$,
если $E_m\le\varepsilon$; везде $\varepsilon=0.05$. Лучшим считается
метод с наименьшей ошибкой $E_m$. Лучший метод не обязательно достигает
точности $\varepsilon$: ошибки обоих сравниваемых методов могут быть
больше $0.05$.

\subsubsection{Набор уравнений}
Уравнения строятся методом подставленного решения: сначала задаётся
точное решение $f$ --- гладкая функция, равная нулю вне некоторого
отрезка, затем правая часть $g$ вычисляется подстановкой $f$
в уравнение~\eqref{eq:ode}. Коэффициенты задаются формулами
\begin{equation}\label{eq:ml-coef}
\begin{gathered}
a=A_{\mathrm{mid}}+A_{\mathrm{half}}\tanh T_a(x),\qquad
b=B\tanh T_b(x),\\
c=C\tanh T_c(x),\qquad
\lambda=C+\delta,
\end{gathered}
\end{equation}
где $T_a,T_b,T_c$ --- формулы, составленные из элементарных функций
(виды операций перечислены в приложении, раздел П.4),
$0<A_{\mathrm{half}}<A_{\mathrm{mid}}$ и $\delta>0$. Поскольку
$|\tanh u|<1$, при всех $x$ выполняются оценки
\[
0<A_-\le a(x)\le A_+,\qquad |b(x)|\le B,\qquad |c(x)|\le C,\qquad
\lambda-C=\delta>0,
\]
где $A_\pm=A_{\mathrm{mid}}\pm A_{\mathrm{half}}$. Тем самым коэффициент
$a$ положителен и $\lambda>\sup_x|c(x)|$ на всей прямой.

Набор содержит 1000 уравнений: 200 семейств по 5 уравнений. Семейство ---
это формулы одного вида с разными числовыми параметрами. Уравнения
разделены на четыре уровня сложности L1--L4 (по 250 уравнений) и на пять
режимов поведения коэффициентов (по 200 уравнений): постоянные, медленно
меняющиеся, умеренно меняющиеся, осциллирующие и с большим разбросом
значений.

Для трёх уравнений расчёт методом Ремизова завершился аварийно,
и результат не сохранился. Эти уравнения исключены, поэтому сравнение
методов Ремизова и Монте-Карло и обучение модели ведутся
на 997 уравнениях. Отсутствие результата не означает, что метод к этим
уравнениям неприменим. Исключены уравнения
\texttt{AMV1\_F0134\_S03}, \texttt{AMV1\_F0136\_S01}
и \texttt{AMV1\_F0139\_S00}.

\subsubsection{Численные методы}\label{sec:ml-methods}
\paragraph{Метод Ремизова\\}
Метод использует оператор
\[
(S(t)\phi)(x)=\tfrac14\phi\bigl(x+2\sqrt{a(x)t}\bigr)
+\tfrac14\phi\bigl(x-2\sqrt{a(x)t}\bigr)
+\tfrac12\phi\bigl(x+2b(x)t\bigr)+t\,c(x)\phi(x).
\]
Приближённое решение при числе композиций $n$ задаётся формулой
\begin{equation}\label{eq:fn}
f_n(x)=\int_0^\infty e^{-\lambda t}\bigl(S(t/n)^ng\bigr)(x)\,\dd t.
\end{equation}
Для вычисления формулы~\eqref{eq:fn} используется следующий алгоритм:
\begin{enumerate}
\item Выбирается окно $[L,R]$, более широкое, чем отрезок $[-1,1]$,
и строится равномерная сетка. На ней вычисляются коэффициенты
и значения $g$.
\item Интеграл по $t$ обрезается на отрезке $[0,T]$, где
\[
T=\max\left(0,\ \frac1\delta\ln\frac{2}{\delta\,\varepsilon_T}\right),
\qquad \delta=\lambda-C,
\]
и выбирается сетка узлов $t_j$ квадратурной формулы по $t$.
\item Для каждого узла $t_j$ задаётся $\tau_j=t_j/n$ и $u^{(0)}=g$;
последовательно вычисляются $u^{(k+1)}=S(\tau_j)u^{(k)}$,
$k=0,\dots,n-1$. Значения в сдвинутых точках, не совпадающих с узлами
сетки, восстанавливаются кубическим сплайном; за краем окна функция
плавно гасится до нуля.
\item Полученное $u^{(n)}$ умножается на $e^{-\lambda t_j}$, и интеграл
по $t$ приближается по формуле Симпсона.
\end{enumerate}

\paragraph{Метод Монте-Карло\\}
Решение записывается как математическое ожидание (формула
Фейнмана--Каца):
\[
f(x_0)=\mathbb E\int_0^\infty
e^{-\lambda t+\int_0^t c(X_s)\,\dd s}\,g(X_t)\,\dd t,
\]
где $X_t$ --- случайный процесс, заданный уравнением
\[
\dd X=b(X)\,\dd t+\sqrt{2a(X)}\,\dd W,\qquad X_0=x_0,
\]
а $W$ --- винеровский процесс.
Алгоритм:
\begin{enumerate}
\item Интеграл по $t$ обрезается на конечном отрезке $[0,T]$.
\item Из точки $x_0$ запускаются $M$ траекторий по схеме
Эйлера--Маруямы:
\[
X_{k+1}=X_k+b(X_k)\,\Delta t+\sqrt{2a(X_k)\,\Delta t}\;\xi_k,
\]
где $\xi_k$ --- независимые случайные величины с распределением
$\mathcal N(0,1)$.
\item Вдоль каждой траектории интеграл по $t$ вычисляется по формуле
трапеций; получается число $S_j$, $j=1,\dots,M$.
\item $f(x_0)\approx\frac1M\sum_{j=1}^MS_j$.
\end{enumerate}
В работе $x_0=0$. Выполнены четыре независимых расчёта
по $M=1024$ траектории, и используется их среднее:
\[
E_{\mathrm{MC}}=\Bigl|\frac14\sum_{s=1}^4 f_{\mathrm{MC},s}(0)-f(0)\Bigr|.
\]
Это ошибка среднего значения, а не средняя ошибка отдельных расчётов.

\paragraph{Метод Рунге--Кутты (RK4)\\}
Уравнение~\eqref{eq:ode} записывается как система для $Y=(f,v)$, где
$v=f'$:
\[
Y'=F(x,Y),\qquad
F(x,Y)=\Bigl(v,\ \frac{(\lambda-c(x))f-b(x)v-g(x)}{a(x)}\Bigr).
\]
На каждом шаге $h$ вычисляются
\begin{enumerate}
\item $k_1=F(x_k,\ Y_k)$,
\item $k_2=F\bigl(x_k+\tfrac h2,\ Y_k+\tfrac h2k_1\bigr)$,
\item $k_3=F\bigl(x_k+\tfrac h2,\ Y_k+\tfrac h2k_2\bigr)$,
\item $k_4=F(x_k+h,\ Y_k+hk_3)$,
\item $Y_{k+1}=Y_k+\tfrac h6\,(k_1+2k_2+2k_3+k_4)$.
\end{enumerate}
Начальное значение $Y_0=\bigl(f(-1),f'(-1)\bigr)$.

\paragraph{Метод мультипликативного интеграла\\}
Уравнение записывается как система для вектора $Z=(f,f',1)^T$:
\[
Z'=B(x)\,Z,\qquad
B(x)=\begin{pmatrix}
0 & 1 & 0\\[2pt]
\dfrac{\lambda-c}{a} & -\dfrac{b}{a} & -\dfrac{g}{a}\\[8pt]
0 & 0 & 0
\end{pmatrix}.
\]
Шаг делается по формуле
\[
Z_{k+1}=\exp\bigl(hB(x_k)\bigr)Z_k,
\]
то есть на каждом шаге коэффициенты и правая часть берутся в левой точке
и считаются постоянными. Матричная экспонента вычисляется точно,
но ошибка от такой замены переменных коэффициентов постоянными остаётся.

\paragraph{Настройки методов\\}
Настройки фиксированы и одинаковы для всех уравнений
(таблица~\ref{tab:ml-settings}); для отдельных уравнений они
не подбирались.

\begin{table}[!htbp]
\centering
\caption{Настройки методов.}
\label{tab:ml-settings}
\small
\begin{tabular}{>{\raggedright\arraybackslash}p{0.3\textwidth}>{\raggedright\arraybackslash}p{0.62\textwidth}}
\toprule
Метод & Настройки\\
\midrule
Метод Ремизова & $n=128$, $\varepsilon_T=10^{-8}$; 513 узлов квадратуры по $t$; шаг сетки по $x$ равен $0.0125$; кубический сплайн и плавное затухание за краем окна\\
Монте-Карло & 4 расчёта по 1024 траектории; шаг схемы Эйлера--Маруямы $\Delta t=1/128$; используется среднее четырёх значений в точке $x=0$\\
RK4 & 128 шагов на $[-1,1]$, $h=1/64$; начальные условия $f(-1)$, $f'(-1)$\\
Мультипликативный интеграл & та же сетка и те же начальные условия, что у RK4\\
\bottomrule
\end{tabular}
\end{table}

Сетка RK4 и метода мультипликативного интеграла содержит 129 узлов,
среди которых есть узел $x=0$, поэтому значение в нуле берётся прямо
в узле сетки. Шаги разных методов имеют разный смысл, поэтому объём
вычислений у методов не одинаков.

\subsubsection{Признаки уравнения}
Модель получает на вход только величины, известные до решения
уравнения: значения функций $a,b,c,g$ на равномерной сетке
$x_i=-1+2i/255$, $i=0,\dots,255$, числа $A_-,A_+,B,C,\lambda$
из~\eqref{eq:ml-coef} и формулы $T_a,T_b,T_c$. По ним строятся
258 признаков (полный перечень приведён в приложении):
\begin{itemize}
\item 128 значений: по 32 значения каждой из функций $a,b,c,g$
в почти равноотстоящих узлах сетки;
\item 56 статистик: для каждой из четырёх функций 14 величин
(минимум, максимум, среднее, разброс, оценки первой и второй
производных, число смен знака и другие);
\item 6 коэффициентов корреляции между парами функций;
\item 48 характеристик формул $T_a,T_b,T_c$: число операций каждого вида
и глубина вложения;
\item 5 границ коэффициентов: $A_-$, $A_+$, $B$, $C$ и число $\lambda$;
\item 10 величин, вычисляемых по этим границам: $\delta=\lambda-C$,
верхний предел $T$ интеграла по $t$ в методе Ремизова, наибольшие сдвиги
аргумента в операторе $S(\tau)$ и другие;
\item 5 безразмерных отношений, сравнивающих величины слагаемых
уравнения~\eqref{eq:ode}.
\end{itemize}
Точное решение, его производные и начальные условия, ошибки методов
и время счёта в признаки не входят.

\subsubsection{Модель выбора из двух методов и критерии качества}
Каждому из 997 уравнений ставится метка
\[
y=\begin{cases}
1, & E_{\mathrm{MC}}<E_{\mathrm{R}}\quad(\text{точнее Монте-Карло}),\\
0, & \text{иначе (точнее метод Ремизова)},
\end{cases}
\]
где $E_{\mathrm{R}}$ и $E_{\mathrm{MC}}$ --- ошибки~\eqref{eq:ml-error}
методов Ремизова и Монте-Карло. Метод Ремизова точнее на 881 уравнении,
Монте-Карло --- на 116.

Набор разбит по семействам на обучающую (Train, 747 уравнений),
валидационную (Validation, 100) и тестовую (Test, 150) выборки;
Монте-Карло точнее в 94, 13 и 9 уравнениях этих выборок соответственно.
Все пять уравнений одного семейства попадают в одну выборку, поэтому
на Validation и Test модель проверяется на формулах нового вида,
а не на тех же формулах с другими параметрами. Тестовая выборка ранее
использовалась в других расчётах проекта, поэтому оценки на ней нельзя
считать полностью независимой проверкой.

Используется модель CatBoost~--- сумма 120 решающих деревьев глубины~2.
Каждое дерево по двум сравнениям вида <<признак меньше порога>> относит
уравнение к одной из четырёх групп и добавляет к сумме число, своё
для каждой группы. По этой сумме модель вычисляет число $p\in(0,1)$.
Монте-Карло выбирается, если $p>t$, иначе выбирается метод Ремизова;
порог $t=0.201$ подобран на валидационной выборке. Параметры обучения:
логистическая функция потерь, шаг обучения $0.05$, коэффициент
$L_2$-регуляризации 10, одинаковые веса классов. Число $p$ не является
оценкой точности численного решения.

Для оценки модели используются следующие величины. Пусть TP --- число
уравнений, где модель выбрала Монте-Карло и он действительно точнее;
FP --- выбран Монте-Карло, но точнее метод Ремизова; FN --- выбран метод
Ремизова, но точнее Монте-Карло. Тогда
\[
P_{\mathrm{MC}}=\frac{TP}{TP+FP},\qquad R_{\mathrm{MC}}=\frac{TP}{TP+FN}
\]
--- доля верных среди выборов Монте-Карло и доля найденных уравнений,
где он точнее. Accuracy --- доля всех верных выборов. Для каждого
из двух классов вычисляется $F_1=2PR/(P+R)$, где $P$ и $R$ определяются
так же, как $P_{\mathrm{MC}}$ и $R_{\mathrm{MC}}$; macro-F1 --- среднее
этих двух значений. Классы сильно несбалансированы, поэтому accuracy
сама по себе малоинформативна: на тестовой выборке постоянный выбор
метода Ремизова уже даёт accuracy 94\% (141 из 150).

Основной критерий --- средняя избыточная ошибка
\begin{equation}\label{eq:ml-regret}
r_i=E_{\text{выбр},i}-\min\bigl(E_{\mathrm{R},i},\,E_{\mathrm{MC},i}\bigr)\ge0,\qquad
\bar r=\frac1N\sum_{i=1}^N r_i,
\end{equation}
где $E_{\text{выбр},i}$ --- ошибка выбранного метода на $i$-м уравнении.
Величина $\bar r$ показывает, на сколько в среднем ошибка выбранного
метода больше ошибки лучшего из двух. Сумма
$\sum_i\min(E_{\mathrm{R},i},E_{\mathrm{MC},i})$ от модели не зависит,
поэтому минимизация $\bar r$ равносильна минимизации средней ошибки
выбранного метода. Дополнительно подсчитывается число уравнений,
на которых выбранный метод достигает точности $0.05$.

\subsubsection{Модель выбора из трёх методов}\label{sec:ml3-setup}
Отдельно рассмотрена задача выбора из трёх методов: метода Ремизова
($\mathrm{R}$), метода Монте-Карло ($\mathrm{MC}$) и метода
мультипликативного интеграла ($\mathrm{PI}$). Для каждого уравнения
лучшим считается метод с наименьшей ошибкой~\eqref{eq:ml-error}:
\[
m_*=\arg\min_{m\in\{\mathrm{R},\,\mathrm{MC},\,\mathrm{PI}\}}E_m .
\]
Достижение точности $E_m\le0.05$ проверяется отдельно и на определение
лучшего метода не влияет. Используются уже полученные ошибки методов
с настройками из таблицы~\ref{tab:ml-settings}; новых расчётов
не проводилось. Лучший метод определён для тех же 997 уравнений: для трёх
уравнений нет результата метода Ремизова.

Выборки Train, Validation и Test те же, что выше, и разбиты
по семействам. Число уравнений, на которых точнее каждый из трёх
методов, приведено в таблице~\ref{tab:ml3-split}.

\begin{table}[!htbp]
\centering
\caption{Число уравнений, на которых каждый из трёх методов даёт
наименьшую ошибку в точке $x=0$.}
\label{tab:ml3-split}
\resizebox{\textwidth}{!}{%
\begin{tabular}{lccccc}
\toprule
Выборка & Уравнений & Семейств & \begin{tabular}[b]{@{}c@{}}Метод\\Ремизова\end{tabular} & \begin{tabular}[b]{@{}c@{}}Монте-\\Карло\end{tabular} & \begin{tabular}[b]{@{}c@{}}Мультипликативный\\интеграл\end{tabular}\\
\midrule
Train & 747 & 150 & 608 & 87 & 52\\
Validation & 100 & 20 & 82 & 12 & 6\\
Test & 150 & 30 & 136 & 7 & 7\\
\midrule
Всего & 997 & 200 & 826 & 106 & 65\\
\bottomrule
\end{tabular}}
\end{table}

Используется модель CatBoost с теми же 258 признаками --- сумма
240 решающих деревьев глубины 3. Каждое дерево по трём сравнениям вида
<<признак меньше порога>> относит уравнение к одной из восьми групп.
Шаг обучения равен $0.05$, коэффициент $L_2$-регуляризации --- 10.
По признакам уравнения модель указывает один из трёх методов.

Прежде чем выбрать эту модель, проверялись разные наборы признаков,
веса классов и параметры модели, а также варианты, в которых модель
предсказывает не сам метод, а величины, связанные с ошибками методов
(регрессия). Всего выполнено 320 обучений: 120 основных и 200
при перекрёстной проверке по семействам. При перекрёстной проверке
набор несколько раз делится по семействам на части; модель обучается
на всех частях, кроме одной, и проверяется на оставшейся. Модель
и правило выбора метода выбраны по валидационной выборке и зафиксированы
до проверки на тестовой.

Качество оценивается следующими величинами: accuracy --- доля уравнений,
на которых выбран метод с наименьшей ошибкой; macro-F1 --- среднее
значений $F_1$ по трём методам (для каждого метода $F_1$ вычисляется так
же, как в предыдущем разделе); число уравнений, на которых выбранный
метод достигает точности $0.05$; средняя ошибка выбранного метода.
Кроме того, строится матрица ошибок: её элемент в строке $i$ и столбце
$j$ --- число уравнений, на которых точнее всего метод $i$, а модель
выбрала метод $j$.

Программы, набор уравнений и сохранённые результаты расчётов доступны
в репозитории~\cite{repo2488}.

\subsection{Результаты}
\subsubsection{Сравнение методов Ремизова и Монте-Карло}
На 997 уравнениях метод Ремизова точнее в 881 случае, Монте-Карло ---
в 116 (рис.~\ref{fig:ml-wins2}). Медиана ошибки метода Ремизова равна
$0.0012$, Монте-Карло --- $0.0092$. Однако у метода Ремизова встречаются
редкие очень большие ошибки (наибольшая равна 842), тогда как ошибка
Монте-Карло не превышает $0.57$ (рис.~\ref{fig:ml-errors}). Точности
$0.05$ метод Ремизова достигает на 982 уравнениях из 997,
Монте-Карло --- на 935.

\begin{figure}[!htbp]
\centering
\includegraphics[width=\textwidth]{simple_01_two_methods_wins.png}
\caption{Число уравнений, на которых каждый из методов даёт меньшую
ошибку в точке $x=0$ (997 уравнений).}
\label{fig:ml-wins2}
\end{figure}

\begin{figure}[!htbp]
\centering
\includegraphics[width=\textwidth]{01_R_MC_error_comparison.png}
\caption{Ошибки методов Ремизова и Монте-Карло в точке $x=0$
на 997 уравнениях (логарифмические шкалы). Слева: каждая точка ---
одно уравнение; точки ниже диагонали соответствуют уравнениям, где
точнее Монте-Карло; пунктир --- уровень $0.05$. Справа: доля уравнений,
ошибка на которых не больше заданной. Remizov R128 --- метод Ремизова
при $n=128$, MC mean4 --- среднее четырёх расчётов Монте-Карло.}
\label{fig:ml-errors}
\end{figure}

\subsubsection{Модель выбора из двух методов}\label{sec:ml-res2}
Показатели модели на трёх выборках приведены в таблице~\ref{tab:ml-258},
её решения на тестовой выборке --- на рис.~\ref{fig:ml-test}. Показатели
на обучающей выборке вычислены на тех же данных, на которых модель
обучалась, и поэтому завышены.

\begin{table}[!htbp]
\centering
\caption{Показатели модели. $P_{\mathrm{MC}}$ --- доля верных
среди выборов Монте-Карло, $R_{\mathrm{MC}}$ --- доля найденных уравнений,
где точнее Монте-Карло, $\bar r$ --- средняя избыточная
ошибка~\eqref{eq:ml-regret}.}
\label{tab:ml-258}
\resizebox{\textwidth}{!}{%
\begin{tabular}{lccccccc}
\toprule
Выборка & Accuracy & macro-F1 & $P_{\mathrm{MC}}$ & $R_{\mathrm{MC}}$ & средняя $E$ & $\bar r$ & $E\le0.05$\\
\midrule
Train (747) & 91.97\% & 0.803 & 0.718 & 0.596 & 0.00506 & 0.00075 & 733/747\\
Validation (100) & 89.00\% & 0.707 & 0.625 & 0.385 & 0.00347 & 0.00088 & 99/100\\
Test (150) & 88.67\% & 0.676 & 0.300 & 0.667 & 0.00560 & 0.00171 & 145/150\\
\bottomrule
\end{tabular}}
\end{table}

\begin{figure}[!htbp]
\centering
\includegraphics[width=\textwidth]{simple_02_two_methods_model_test.png}
\caption{Решения модели на тестовой выборке (150 уравнений). Верным
считается выбор метода с меньшей ошибкой; это не гарантирует достижения
точности $0.05$.}
\label{fig:ml-test}
\end{figure}

В таблице~\ref{tab:ml-rules} модель сравнивается с простыми правилами
выбора на тестовой выборке. Идеальный выбор --- выбор по известным
ошибкам; до решения уравнения он недоступен и показывает предел,
к которому можно стремиться.

\begin{table}[!htbp]
\centering
\caption{Сравнение модели с простыми правилами выбора на тестовой
выборке (150 уравнений).}
\label{tab:ml-rules}
\begin{tabular}{lcccc}
\toprule
Правило выбора & средняя $E$ & $\bar r$ & TP/FP/FN & $E\le0.05$\\
\midrule
модель & 0.00560 & 0.00171 & 6/14/3 & 145/150\\
всегда метод Ремизова & 1.659 & 1.655 & 0/0/9 & 147/150\\
всегда Монте-Карло & 0.01728 & 0.01338 & 9/141/0 & 141/150\\
идеальный выбор & 0.00390 & 0 & 9/0/0 & 147/150\\
\bottomrule
\end{tabular}
\end{table}

Модель находит 6 из 9 уравнений, где точнее Монте-Карло, но лишь 30\%
её выборов Монте-Карло верны: из 20 таких выборов ошибочны 14. Число
уравнений, на которых достигнута точность $0.05$, уменьшается со 147
при постоянном выборе метода Ремизова до 145.

Большая разница средних ошибок модели и постоянного выбора метода
Ремизова требует оговорки. Она почти целиком определяется одним
уравнением \texttt{AMV1\_F0180\_S04}, где $E_{\mathrm{R}}=248.3$,
а $E_{\mathrm{MC}}=0.103$: при выборе метода Ремизова избыточная ошибка
на нём равна $248.2$, что после деления на 150 даёт $1.655$. Модель верно
выбирает на этом уравнении Монте-Карло, но точность $0.05$ всё равно
не достигнута. Если для проверки устойчивости вывода исключить это
уравнение, то на остальных 149 средняя ошибка постоянного выбора метода
Ремизова равна $0.00350$, а модели --- $0.00495$, то есть преимущество
меняет знак. Во всех основных результатах это уравнение сохранено.
Насколько результат переносится на другие семейства уравнений,
не установлено.

\subsubsection{Разбор неверных выборов}
На тестовой выборке модель ошиблась 17 раз: 14 раз выбрала Монте-Карло,
когда точнее метод Ремизова (FP), и 3 раза выбрала метод Ремизова, когда
точнее Монте-Карло (FN). Все случаи приведены
в таблице~\ref{tab:ml-wrong} и на рис.~\ref{fig:ml-wrong}; для каждого
из них избыточная ошибка равна $|\Delta|=|E_{\mathrm{R}}-E_{\mathrm{MC}}|$.

\begin{table}[!htbp]
\centering
\caption{Неверные выборы модели на тестовой выборке.
$\Delta=E_{\mathrm{R}}-E_{\mathrm{MC}}$; отрицательное значение означает,
что точнее метод Ремизова. Значения округлены.}
\label{tab:ml-wrong}
\resizebox{\textwidth}{!}{%
\begin{tabular}{llcccl}
\toprule
Уравнение & Тип & $E_{\mathrm{R}}$ & $E_{\mathrm{MC}}$ & $\Delta$ & Точность $0.05$\\
\midrule
\texttt{AMV1\_F0036\_S03} & FP & 0.001272 & 0.003581 & $-0.002310$ & у обоих методов\\
\texttt{AMV1\_F0084\_S00} & FP & 0.005260 & 0.01419 & $-0.008927$ & у обоих методов\\
\texttt{AMV1\_F0084\_S03} & FP & 0.005535 & 0.01253 & $-0.006993$ & у обоих методов\\
\texttt{AMV1\_F0094\_S01} & FP & 0.0007748 & 0.02212 & $-0.02134$ & у обоих методов\\
\texttt{AMV1\_F0098\_S02} & FP & 0.0008599 & 0.003362 & $-0.002502$ & у обоих методов\\
\texttt{AMV1\_F0098\_S04} & FP & 0.001981 & 0.01702 & $-0.01504$ & у обоих методов\\
\texttt{AMV1\_F0100\_S04} & FP & 0.0004919 & 0.006008 & $-0.005516$ & у обоих методов\\
\texttt{AMV1\_F0164\_S01} & FP & 0.0002026 & 0.004107 & $-0.003904$ & у обоих методов\\
\texttt{AMV1\_F0179\_S03} & FP & 0.008291 & 0.01010 & $-0.001805$ & у обоих методов\\
\texttt{AMV1\_F0180\_S01} & FP & 0.03175 & 0.09655 & $-0.06480$ & \textbf{теряется}\\
\texttt{AMV1\_F0180\_S02} & FP & 0.003379 & 0.06041 & $-0.05704$ & \textbf{теряется}\\
\texttt{AMV1\_F0190\_S01} & FP & 0.009266 & 0.04533 & $-0.03606$ & у обоих методов\\
\texttt{AMV1\_F0191\_S03} & FP & 0.001227 & 0.002277 & $-0.001050$ & у обоих методов\\
\texttt{AMV1\_F0191\_S04} & FP & 0.01014 & 0.03150 & $-0.02136$ & у обоих методов\\
\texttt{AMV1\_F0070\_S02} & FN & 0.002124 & 0.001312 & $0.0008121$ & у обоих методов\\
\texttt{AMV1\_F0190\_S04} & FN & 0.003905 & 0.0004051 & $0.003500$ & у обоих методов\\
\texttt{AMV1\_F0191\_S01} & FN & 0.005347 & 0.002114 & $0.003233$ & у обоих методов\\
\bottomrule
\end{tabular}}
\end{table}

\begin{figure}[!htbp]
\centering
\includegraphics[width=0.85\textwidth]{simple_03_wrong_mc_choices.png}
\caption{Избыточная ошибка $|E_{\mathrm{R}}-E_{\mathrm{MC}}|$
в 17 неверных выборах модели; цвет показывает тип ошибки.}
\label{fig:ml-wrong}
\end{figure}

В 12 из 14 случаев FP оба метода достигают точности $0.05$, но разница
ошибок не всегда мала. Например, для \texttt{AMV1\_F0190\_S01} ошибка
возрастает с $0.0093$ до $0.0453$, почти до порога. В двух случаях
неверный выбор приводит к потере точности: для \texttt{AMV1\_F0180\_S01}
ошибка возрастает с $0.0317$ до $0.0965$, для \texttt{AMV1\_F0180\_S02}
--- с $0.0034$ до $0.0604$; избыточные ошибки равны $0.0648$ и $0.0570$.

Во всех трёх случаях FN оба метода достигают точности $0.05$;
упущенное уменьшение ошибки лежит между $0.0008$ и $0.0035$. Суммарная
избыточная ошибка 14 случаев FP равна $0.2486$, трёх случаев FN ---
$0.0075$; вместе $0.2562$, что даёт $\bar r=0.00171$ на 150 уравнениях.
Ни в одном случае модель не заменила метод Ремизова, не достигающий
точности, на Монте-Карло, достигающий её. На двух уравнениях, где модель
верно выбрала Монте-Карло, \texttt{AMV1\_F0070\_S04}
и \texttt{AMV1\_F0180\_S04}, ошибки равны $0.0563$ и $0.1031$, то есть
обе больше $0.05$. Таким образом, умение выбрать более точный метод
и достижение требуемой точности --- разные свойства.

\subsubsection{Сравнение четырёх методов}
К результатам методов Ремизова и Монте-Карло добавлены результаты одного
расчёта методами RK4 и мультипликативного интеграла с настройками
из таблицы~\ref{tab:ml-settings}. Ошибки в точке $x=0$ на общих
997 уравнениях приведены в таблице~\ref{tab:ml-four}.

\begin{table}[!htbp]
\centering
\caption{Ошибки $E$ в точке $x=0$ на 997 уравнениях. Квантиль $0.95$ ---
значение, которое не превышают ошибки 95\% уравнений.}
\label{tab:ml-four}
\resizebox{\textwidth}{!}{%
\begin{tabular}{>{\raggedright\arraybackslash}p{4.8cm}ccccc}
\toprule
Метод & средняя $E$ & медиана & \begin{tabular}[b]{@{}c@{}}квантиль\\$0.95$\end{tabular} & максимум & $E\le0.05$\\
\midrule
Метод Ремизова & 1.127 & 0.00119 & 0.0135 & 842.1 & 982/997\\
Монте-Карло & 0.0181 & 0.00919 & 0.0583 & 0.570 & 935/997\\
RK4 & $1.79\cdot10^{-7}$ & $6.24\cdot10^{-9}$ & $3.89\cdot10^{-7}$ & $4.59\cdot10^{-5}$ & 997/997\\
Мультипликативный интеграл & 0.0456 & 0.0174 & 0.163 & 1.985 & 758/997\\
\bottomrule
\end{tabular}}
\end{table}

На всех 997 уравнениях наименьшую ошибку даёт RK4, совпадений ошибок нет
(рис.~\ref{fig:ml-wins4}); по уровням сложности L1--L4 это 250, 250, 247
и 250 уравнений. Если брать все уравнения, для которых есть результат
метода, то RK4 достигает точности $0.05$ на 1000 из 1000, Монте-Карло ---
на 938 из 1000, метод мультипликативного интеграла --- на 760 из 1000,
метод Ремизова --- на 982 из 997. Знаменатели здесь разные, поэтому
методы попарно сравниваются только по таблице~\ref{tab:ml-four}.

\begin{figure}[!htbp]
\centering
\includegraphics[width=\textwidth]{simple_04_four_methods_wins.png}
\caption{Число уравнений, на которых каждый из четырёх методов даёт
наименьшую ошибку в точке $x=0$ (997 уравнений).}
\label{fig:ml-wins4}
\end{figure}

Средние ошибки показаны на рис.~\ref{fig:ml-mean4}. Редкие очень большие
ошибки метода Ремизова сильно увеличивают его среднее, поэтому
при сравнении нужно смотреть и на медиану в таблице~\ref{tab:ml-four}.
Число уравнений, на которых каждый метод достигает точности $0.05$,
показано на рис.~\ref{fig:ml-success4}.

\begin{figure}[!htbp]
\centering
\includegraphics[width=\textwidth]{simple_05_four_methods_mean_error.png}
\caption{Средняя ошибка четырёх методов в точке $x=0$ на 997 уравнениях.
Шкала линейная, поэтому столбец RK4 почти не виден.}
\label{fig:ml-mean4}
\end{figure}

\begin{figure}[!htbp]
\centering
\includegraphics[width=\textwidth]{simple_06_four_methods_success_rate.png}
\caption{Число уравнений, на которых метод достигает точности $0.05$
в точке $x=0$ (997 уравнений). Достижение точности и наименьшая ошибка ---
разные свойства.}
\label{fig:ml-success4}
\end{figure}

Результат относится к выбранным настройкам методов и к точке $x=0$.
Например, метод мультипликативного интеграла достигает точности $0.05$
в нуле на 760 уравнениях, а по наибольшей ошибке на 33 узлах отрезка
$[-1,1]$ --- лишь на 404. Расчёт с одним шагом не позволяет оценить
порядок сходимости методов, их устойчивость при других условиях и объём
вычислений, нужный для достижения точности. Время счёта методов в этой
серии не сравнивалось.

Для машинного обучения вывод однозначен: на всех 997 уравнениях лучший
метод --- RK4, и постоянный выбор RK4 уже полностью совпадает с метками.
На таком наборе нельзя содержательно проверить, умеет ли модель различать
области, где лучше разные методы, поэтому модель выбора из четырёх
методов не обучалась.

\subsubsection{Выбор из трёх методов}\label{sec:ml3-res}
Если не рассматривать RK4, лучшими оказываются разные методы:
на 997 уравнениях метод Ремизова точнее двух других на 826, Монте-Карло ---
на 106, мультипликативный интеграл --- на 65 (рис.~\ref{fig:ml3-wins},
таблица~\ref{tab:ml3-split}). Результаты модели на тестовой выборке
приведены в таблице~\ref{tab:ml3-test}.

\begin{figure}[!htbp]
\centering
\includegraphics[width=0.8\textwidth]{01_winners.png}
\caption{Число уравнений, на которых каждый из трёх методов даёт
наименьшую ошибку в точке $x=0$ (997 уравнений). R128 --- метод Ремизова
при $n=128$, MC mean4 --- среднее четырёх расчётов Монте-Карло,
PI --- метод мультипликативного интеграла.}
\label{fig:ml3-wins}
\end{figure}

\begin{table}[!htbp]
\centering
\caption{Модель выбора из трёх методов и постоянный выбор одного метода
на тестовой выборке (150 уравнений).}
\label{tab:ml3-test}
\resizebox{\textwidth}{!}{%
\begin{tabular}{lcccc}
\toprule
Правило выбора & Accuracy & macro-F1 & $E\le0.05$ & средняя $E$\\
\midrule
модель & 88.67\% & 0.314 & 147/150 & 1.65885\\
всегда метод Ремизова & 90.67\% & 0.317 & 147/150 & 1.65896\\
всегда Монте-Карло & 4.67\% & 0.030 & 141/150 & 0.01728\\
всегда мультипликативный интеграл & 4.67\% & 0.030 & 110/150 & 0.05416\\
\bottomrule
\end{tabular}}
\end{table}

Модель верно определила лучший метод на 133 из 150 тестовых уравнений
(88.67\%), а выбранный ею метод достиг точности $0.05$ на 147 уравнениях
(98\%). Постоянный выбор метода Ремизова верен на 136 уравнениях
(90.67\%) и тоже достигает точности на 147. Модель выбрала метод
Ремизова 146 раз, Монте-Карло --- 4 раза, мультипликативный интеграл
не выбрала ни разу.

Матрица ошибок модели (строки --- метод, который на самом деле точнее,
столбцы --- выбор модели; порядок: метод Ремизова, Монте-Карло,
мультипликативный интеграл) имеет вид
\[
\begin{pmatrix}
133&3&0\\
7&0&0\\
6&1&0
\end{pmatrix}
\]
(рис.~\ref{fig:ml3-conf}). Все четыре выбора Монте-Карло ошибочны:
в трёх случаях точнее был метод Ремизова, в одном --- мультипликативный
интеграл. Ни одно из 7 уравнений, где точнее Монте-Карло, и ни одно
из 7, где точнее мультипликативный интеграл, модель не нашла. Поэтому
для этих двух методов $F_1=0$, и macro-F1 у модели и у постоянного
выбора метода Ремизова почти одинаков и мал (рис.~\ref{fig:ml3-acc}).

\begin{figure}[!htbp]
\centering
\begin{minipage}[t]{0.49\textwidth}
\centering
\includegraphics[width=\linewidth]{02_accuracy.png}
\end{minipage}\hfill
\begin{minipage}[t]{0.49\textwidth}
\centering
\includegraphics[width=\linewidth]{03_macro_f1.png}
\end{minipage}
\caption{Доля верных выборов (слева) и macro-F1 (справа) на тестовой
выборке для модели и для постоянного выбора метода Ремизова
(на рисунке --- <<Всегда R128>>).}
\label{fig:ml3-acc}
\end{figure}

\begin{figure}[!htbp]
\centering
\includegraphics[width=0.6\textwidth]{05_confusion_matrix.png}
\caption{Матрица ошибок модели выбора из трёх методов на тестовой
выборке: по строкам --- метод, который на самом деле точнее,
по столбцам --- метод, выбранный моделью. R128 --- метод Ремизова,
MC mean4 --- Монте-Карло, PI --- мультипликативный интеграл.}
\label{fig:ml3-conf}
\end{figure}

Средние ошибки модели и постоянного выбора метода Ремизова почти
совпадают (рис.~\ref{fig:ml3-mean}). Значение $1.659$ для постоянного
выбора метода Ремизова то же, что в таблице~\ref{tab:ml-rules}, и, как
показано в разделе~\ref{sec:ml-res2}, почти целиком определяется одним
уравнением \texttt{AMV1\_F0180\_S04}, где $E_{\mathrm{R}}=248.3$. Средняя
ошибка модели почти так же велика; значит, на этом уравнении модель
тоже выбрала метод Ремизова.

\begin{figure}[!htbp]
\centering
\includegraphics[width=0.85\textwidth]{04_mean_error.png}
\caption{Средняя ошибка выбранного метода в точке $x=0$ на тестовой
выборке для модели и для постоянного выбора одного метода. R128 ---
метод Ремизова, MC mean4 --- Монте-Карло, PI --- мультипликативный
интеграл.}
\label{fig:ml3-mean}
\end{figure}

На рис.~\ref{fig:ml3-imp} показаны пять признаков с наибольшей
важностью. Важность --- оценка CatBoost того, насколько признак в среднем
влияет на ответы модели; сумма важностей всех 258 признаков равна 100\%.
Наиболее важны среднеквадратичное значение $b$ (10.36\%; приложение,
П.2, величина 5), значение $b$ в узле $x\approx0.545$ (5.15\%; П.1),
число $\lambda$ (4.85\%; П.5), изменчивость коэффициента $a$
(4.62\%; П.7, признак 3) и величина $\sqrt{2A_+T}$ (4.27\%; П.6,
величина 7). Вместе на них приходится около 29\% суммарной важности.
Важность показывает, какими признаками модель пользуется больше всего,
но не объясняет, почему один метод точнее другого.

\begin{figure}[!htbp]
\centering
\includegraphics[width=0.85\textwidth]{06_top5_features.png}
\caption{Пять признаков с наибольшей важностью в модели выбора из трёх
методов. <<Значение b (b\_sample\_24)>> --- значение $b$ в узле
$x\approx0.545$; <<Диффузионная длина>> --- величина $\sqrt{2A_+T}$.}
\label{fig:ml3-imp}
\end{figure}

Таким образом, построен и проверен полный процесс автоматического выбора
одного из трёх методов по признакам уравнения, но на тестовой выборке
модель не лучше постоянного выбора метода Ремизова: доля верных выборов
у неё на 2 процентных пункта меньше, а число уравнений с достигнутой
точностью и средняя ошибка практически те же. Результат относится
к фиксированным настройкам методов и к ошибке в точке $x=0$; тестовая
выборка ранее использовалась в других расчётах проекта, поэтому оценка
не является полностью независимой.

\section{Заключение}
\subsection{Краткие выводы по полученным результатам}
\begin{enumerate}
\item Построен набор из 1000 уравнений с известными точными решениями,
в котором коэффициент $a$ положителен и $\lambda>\sup_x|c(x)|$
по построению, и получены ошибки четырёх численных методов в точке
$x=0$.
\item Метод Ремизова точнее Монте-Карло на 881 уравнении из 997
и достигает точности $0.05$ на 982 из них, но изредка его ошибка
на несколько порядков больше (до 842).
\item Модель выбора находит 6 из 9 тестовых уравнений, где точнее
Монте-Карло, но делает 14 ошибочных выборов и 3 пропуска; два ошибочных
выбора приводят к потере точности $0.05$. Её выигрыш в средней ошибке
определяется одним уравнением, поэтому устойчивого преимущества перед
постоянным выбором метода Ремизова нет. Выводы ограничены повторным
использованием тестовой выборки и малым числом уравнений, где точнее
Монте-Карло.
\item Выбор более точного метода и достижение точности $0.05$ ---
разные свойства: верно выбранный метод может не достигать точности.
\item При выбранных настройках RK4 точнее остальных методов в точке
$x=0$ на всех 997 уравнениях и достигает точности $0.05$ на всех 1000,
поэтому выбор из четырёх методов на этом наборе тривиален. Это результат
для данного набора и данных настроек, а не доказательство превосходства
RK4 в общем случае.
\item Модель выбора из трёх методов --- Ремизова, Монте-Карло и метода
мультипликативного интеграла --- верно определяет лучший метод
на 88.67\% тестовых уравнений, а выбранный ею метод достигает точности $0.05$
на 98\% из них. Однако постоянный выбор метода Ремизова даёт 90.67\%
и те же 98\%: модель почти всегда выбирает метод Ремизова (146 раз
из 150) и не находит ни одного уравнения, где точнее другой метод.
\end{enumerate}

\subsection{Оценка полноты решения поставленных задач и планируемое развитие}
Выполнены все поставленные задачи: построен набор уравнений, вычислены
ошибки четырёх методов, построены признаки, обучена и проверена модель
выбора между методами Ремизова и Монте-Карло, разобраны все её неверные
выборы на тестовой выборке, проведено сравнение четырёх методов, обучена
и проверена модель выбора из трёх методов.

Дальнейшее исследование должно проверять точность и устойчивость методов
при заранее заданных изменениях задач и настроек (в том числе шага
и объёма вычислений), сохраняя независимую тестовую выборку. Для модели
выбора из нескольких методов нужны задачи, на которых лучшими
оказываются разные методы.

\clearpage
\renewcommand{\refname}{Список использованных источников}
\phantomsection
\addcontentsline{toc}{section}{Список использованных источников}
\begin{thebibliography}{9}
\bibitem{repo2488} Timofeev, M. Yu.. ode-generator-remizov: генератор уравнений, решатели и эксперименты по выбору метода. GitHub. 2026. URL: \url{https://github.com/MikTi05/ode-generator-remizov}.

\end{thebibliography}

\clearpage
\section*{Приложение. Признаки модели выбора метода}
\phantomsection
\addcontentsline{toc}{section}{Приложение. Признаки модели выбора метода}

Признаки вычисляются по функциям $v\in\{a,b,c,g\}$ из
уравнения~\eqref{eq:ode} на сетке $x_i=-1+2i/255$, $i=0,\dots,255$;
$v_i=v(x_i)$. Всего признаков
\[
258=128+56+6+48+5+10+5.
\]

\subsection*{П.1. Значения функций (128 признаков)}
Для каждой функции $v\in\{a,b,c,g\}$ берутся 32 значения
\[
v\bigl(x_{j_k}\bigr),\qquad j_k=\lfloor 255k/31\rfloor,\quad k=0,\dots,31,
\]
то есть значения в почти равноотстоящих узлах от $x=-1$ до $x=1$.
Значения берутся в узлах сетки, без интерполяции.

\subsection*{П.2. Статистики функций (56 признаков)}
Для каждой функции $v\in\{a,b,c,g\}$ вычисляются 14 величин из
таблицы~\ref{tab:app-stats}. Обозначения:
\begin{gather*}
\bar v=\frac1{256}\sum_{i=0}^{255}v_i,\qquad
d^{(1)}_i=\frac{v_{i+1}-v_i}{x_{i+1}-x_i},\quad i=0,\dots,254,\\
d^{(2)}_i=\frac{d^{(1)}_{i+1}-d^{(1)}_i}{m_{i+1}-m_i},\quad i=0,\dots,253,
\end{gather*}
где $m_i=(x_i+x_{i+1})/2$ --- середины отрезков сетки;
$d^{(1)}$ и $d^{(2)}$ --- разностные оценки первой и второй производных.
Через $\hat v_k$ обозначено дискретное преобразование Фурье
последовательности $v_i-\bar v$, $P_k=|\hat v_k|^2$, $k=0,\dots,128$.

\begin{table}[!htbp]
\centering
\caption{Статистики, вычисляемые для каждой из функций $a,b,c,g$.}
\label{tab:app-stats}
\renewcommand{\arraystretch}{1.25}
\begin{tabular}{r>{\raggedright\arraybackslash}p{0.4\textwidth}l}
\toprule
№ & Величина & Формула\\
\midrule
1 & максимум & $\max_i v_i$\\
2 & минимум & $\min_i v_i$\\
3 & среднее & $\bar v$\\
4 & стандартное отклонение & $\Bigl(\frac1{256}\sum_i(v_i-\bar v)^2\Bigr)^{1/2}$\\
5 & среднеквадратичное значение & $\Bigl(\frac1{256}\sum_i v_i^2\Bigr)^{1/2}$\\
6 & максимум модуля & $\max_i|v_i|$\\
7 & медиана & медиана чисел $v_0,\dots,v_{255}$\\
8 & полная вариация & $\sum_{i=0}^{254}|v_{i+1}-v_i|$\\
9 & среднеквадратичная первая производная & $\Bigl(\frac1{255}\sum_i\bigl(d^{(1)}_i\bigr)^2\Bigr)^{1/2}$\\
10 & максимум модуля первой производной & $\max_i\bigl|d^{(1)}_i\bigr|$\\
11 & среднеквадратичная вторая производная & $\Bigl(\frac1{254}\sum_i\bigl(d^{(2)}_i\bigr)^2\Bigr)^{1/2}$\\
12 & доля отрицательных значений & $\frac1{256}\,\#\{i:\ v_i<0\}$\\
13 & число смен знака & $\#\{i:\ v_iv_{i+1}<0\}$\\
14 & доля высоких частот & $\sum_{k\ge8}P_k\big/\sum_kP_k$\\
\bottomrule
\end{tabular}
\end{table}
Если $\sum_kP_k=0$ (функция постоянна), величина 14 полагается равной нулю.

\subsection*{П.3. Корреляции (6 признаков)}
Для каждой пары функций $(a,b)$, $(a,c)$, $(a,g)$, $(b,c)$, $(b,g)$,
$(c,g)$ вычисляется коэффициент корреляции значений на сетке
\[
\rho(v,w)=\frac{\sum_i(v_i-\bar v)(w_i-\bar w)}
{\sqrt{\sum_i(v_i-\bar v)^2}\,\sqrt{\sum_i(w_i-\bar w)^2}};
\]
если одна из функций постоянна, полагается $\rho=0$.

\subsection*{П.4. Характеристики формул коэффициентов (48 признаков)}
Формулы $T_a,T_b,T_c$ из~\eqref{eq:ml-coef} записываются в виде дерева
операций (до преобразования $\tanh$). Для каждой из трёх формул
вычисляются 16 величин: общее число узлов дерева, глубина дерева
(корень имеет глубину 1) и число узлов каждого из 14 видов
из таблицы~\ref{tab:app-ops}. Числовые параметры $\alpha,\beta,\omega,
\varphi,\mu$ в признаки не входят. Формула точного решения не используется.

\begin{table}[!htbp]
\centering
\caption{Виды узлов в формулах коэффициентов.}
\label{tab:app-ops}
\renewcommand{\arraystretch}{1.3}
\begin{tabular}{>{\raggedright\arraybackslash}p{0.32\textwidth}l}
\toprule
Тип узла & Вид\\
\midrule
функции от $x$ (листья) & константа;\ $\sin(\omega x+\varphi)$;\ $\cos(\omega x+\varphi)$;\\
 & $\tanh(\omega x+\varphi)$;\ $e^{-\alpha(x-\mu)^2}$;\\
 & $\dfrac{1}{1+\alpha(x-\mu)^2}$;\ $\dfrac{x-\mu}{\sqrt{1+(x-\mu)^2}}$\\
\midrule
функции одного аргумента $u$ & $\sin u$;\ $\cos u$;\ $\tanh u$;\ $e^{-u^2}$;\ $\dfrac{u}{1+u^2}$\\
\midrule
функции двух аргументов $u,w$ & $\alpha u+\beta w$;\ $u\cdot w$\\
\bottomrule
\end{tabular}
\end{table}

\subsection*{П.5. Границы коэффициентов (5 признаков)}
Числа $A_-$, $A_+$, $B$, $C$ и $\lambda$ из~\eqref{eq:ml-coef}:
\[
A_-\le a(x)\le A_+,\qquad |b(x)|\le B,\qquad |c(x)|\le C
\quad\text{при всех } x\in\mathbb R.
\]
Это точные границы, следующие из формул коэффициентов, а не минимумы
и максимумы значений на сетке.

\subsection*{П.6. Величины, вычисляемые по границам (10 признаков)}
Величины из таблицы~\ref{tab:app-derived} вычисляются по границам
из раздела П.5. Здесь $T$ --- верхний предел интеграла по $t$ в методе
Ремизова (раздел~\ref{sec:ml-methods}) при $\varepsilon_T=10^{-8}$,
а $\tau=T/n$ при $n=128$ --- наибольший шаг в операторе $S(\tau)$.

\begin{table}[!htbp]
\centering
\caption{Величины, вычисляемые по границам коэффициентов.}
\label{tab:app-derived}
\renewcommand{\arraystretch}{1.25}
\begin{tabular}{rl>{\raggedright\arraybackslash}p{0.46\textwidth}}
\toprule
№ & Формула & Смысл\\
\midrule
1 & $\delta=\lambda-C$ & запас в условии $\lambda>\sup_x|c(x)|$\\
2 & $T=\max\Bigl(0,\ \dfrac1\delta\ln\dfrac{2}{\delta\cdot10^{-8}}\Bigr)$ & верхний предел интеграла по $t$\\[8pt]
3 & $T/128$ & наибольший шаг $\tau$ в операторе $S(\tau)$\\
4 & $\lambda T$ & скорость затухания множителя $e^{-\lambda t}$ на $[0,T]$\\
5 & $CT$ & произведение границы $|c|$ на $T$\\
6 & $A_+/A_-$ & отношение наибольшего и наименьшего значений $a$\\
7 & $\sqrt{2A_+T}$ & характерный масштаб сдвигов, связанных с $a$, за время $T$\\
8 & $2\sqrt{A_+T/128}$ & наибольший сдвиг $2\sqrt{a(x)\tau}$ в операторе $S(\tau)$\\
9 & $BT$ & характерный масштаб сдвигов, связанных с $b$, за время $T$\\
10 & $2BT/128$ & наибольший сдвиг $2b(x)\tau$ в операторе $S(\tau)$\\
\bottomrule
\end{tabular}
\end{table}

\subsection*{П.7. Отношения слагаемых уравнения (5 признаков)}
Положим $q(x)=\lambda-c(x)>0$; $a_0,b_0,q_0$ --- значения $a,b,q$ в точке
$x=0$, найденные линейной интерполяцией по сетке. Производные вычисляются
разностными формулами второго порядка на той же сетке, интегралы ---
по формуле трапеций на $[-1,1]$; $\operatorname{RMS}(u)=
\bigl(\frac1{256}\sum_iu_i^2\bigr)^{1/2}$.

Если в уравнении~\eqref{eq:ode} пренебречь слагаемыми $af''$ и $bf'$,
то получается простое приближение $f\approx u=g/q$. Признаки сравнивают
величины слагаемых уравнения и показывают, насколько это приближение
нарушается.

\begin{table}[!htbp]
\centering
\caption{Отношения слагаемых уравнения.}
\label{tab:app-extra}
\renewcommand{\arraystretch}{1.6}
\begin{tabular}{rl>{\raggedright\arraybackslash}p{0.46\textwidth}}
\toprule
№ & Формула & Смысл\\
\midrule
1 & $\dfrac{|b_0|}{\sqrt{a_0q_0}}$ & величина коэффициента $b$ по сравнению с $a$ и $q$ в нуле\\[10pt]
2 & $\dfrac{a_0\,\operatorname{RMS}(g'')}{q_0\,\operatorname{RMS}(g)}$ & величина $ag''$ по сравнению с $qg$\\[10pt]
3 & $\dfrac{\operatorname{RMS}(a')}{\sqrt{a_0q_0}}$ & изменчивость коэффициента $a$\\[10pt]
4 & $\dfrac{\bigl|\int_{-1}^{1}g\,dx\bigr|}{\int_{-1}^{1}|g|\,dx}$ & взаимное погашение положительной и отрицательной частей $g$\\[10pt]
5 & $\dfrac{\operatorname{RMS}(au''+bu')}{\operatorname{RMS}(g)},\ u=g/q$ & невязка приближения $f\approx g/q$ в уравнении\\
\bottomrule
\end{tabular}
\end{table}
При $g\equiv0$ признаки 2, 4 и 5 полагаются равными нулю.
\end{document}